%  LaTeX support: latex@mdpi.com 
%  For support, please attach all files needed for compiling as well as the log file, and specify your operating system, LaTeX version, and LaTeX editor.

%=================================================================
\documentclass[sustainability,article,submit,pdftex,moreauthors]{Definitions/mdpi} 

%=================================================================
% MDPI internal commands
\firstpage{1} 
\makeatletter 
\setcounter{page}{\@firstpage} 
\makeatother
\pubvolume{1}
\issuenum{1}
\articlenumber{0}
\pubyear{2022}
\copyrightyear{2022}
%\externaleditor{Academic Editor: Firstname Lastname}
\datereceived{4 November 2022} 
\dateaccepted{} 
\datepublished{} 
%\datecorrected{} % Corrected papers include a "Corrected: XXX" date in the original paper.
%\dateretracted{} % Corrected papers include a "Retracted: XXX" date in the original paper.
\hreflink{https://doi.org/} % If needed use \linebreak
\graphicspath{{figures/}} % path to folder with figures
%\doinum{}

\usepackage{listings}
\usepackage{xcolor}

\definecolor{codegreen}{rgb}{0,0.6,0}
\definecolor{codegray}{rgb}{0.5,0.5,0.5}
\definecolor{codepurple}{RGB}{204, 0, 153}
\definecolor{backcolour}{RGB}{255, 250, 230}
%    
\lstdefinestyle{mystyle}{
    backgroundcolor=\color{backcolour}, 
    commentstyle=\color{codegreen},
    keywordstyle=\color{magenta},
    numberstyle=\tiny\color{codegray},
    stringstyle=\color{codepurple},
    basicstyle=\ttfamily\footnotesize,
    breakatwhitespace=false,         
    breaklines=true,                 
    captionpos=b,                    
    keepspaces=true,                 
    numbers=left,                    
    numbersep=5pt,                  
    showspaces=false,                
    showstringspaces=false,
    showtabs=false,                  
    tabsize=2
}

\lstset{style=mystyle}
	
\usepackage{soul} % highlighting text
%------------------------------------------------------------------
% The following line should be uncommented if the LaTeX file is uploaded to arXiv.org
%\pdfoutput=1

%=================================================================
% Add packages and commands here. The following packages are loaded in our class file: fontenc, inputenc, calc, indentfirst, fancyhdr, graphicx, epstopdf, lastpage, ifthen, lineno, float, amsmath, setspace, enumitem, mathpazo, booktabs, titlesec, etoolbox, tabto, xcolor, soul, multirow, microtype, tikz, totcount, changepage, attrib, upgreek, cleveref, amsthm, hyphenat, natbib, hyperref, footmisc, url, geometry, newfloat, caption

%=================================================================
%% Please use the following mathematics environments: Theorem, Lemma, Corollary, Proposition, Characterization, Property, Problem, Example, ExamplesandDefinitions, Hypothesis, Remark, Definition, Notation, Assumption
%% For proofs, please use the proof environment (the amsthm package is loaded by the MDPI class).

%=================================================================
% Full title of the paper (Capitalized)

\Title{Seismotectonics of shallow-focus earthquakes in Venezuela with links to gravity anomalies and geologic heterogeneity mapped by GMT scripting language}

% MDPI internal command: Title for citation in the left column
\TitleCitation{Seismotectonics of shallow-focus earthquakes in Venezuela with links to gravity anomalies and geologic heterogeneity mapped by GMT scripting language}

% Author Orchid ID: enter ID or remove command
\newcommand{\orcidauthorA}{0000-0002-5759-1089} % Add \orcidA{} behind the author's name
\newcommand{\orcidauthorB}{0000-0002-6461-1551} % Add \orcidB{} behind the author's name

% Authors, for the paper (add full first names)
\Author{Polina Lemenkova $^{1}$*\orcidA{} and Olivier Debeir $^{1}\orcidB{}$}

% MDPI internal command: Authors, for metadata in PDF
\AuthorNames{Poina Lemenkova and Olivier Debeir}

% MDPI internal command: Authors, for citation in the left column
\AuthorCitation{Lemenkova, P.; Debeir, O.}
% If this is a Chicago style journal: Lastname, Firstname, Firstname Lastname, and Firstname Lastname.

% Affiliations / Addresses (Add [1] after \address if there is only one affiliation.)
\address{%
$^{1}$ \quad Université Libre de Bruxelles (ULB), École polytechnique de Bruxelles (Brussels Faculty of Engineering), Laboratory of Image Synthesis and Analysis (LISA). Building L, Campus de Solbosch, ULB -- LISA CP165/57, Avenue Franklin D. Roosevelt 50, Brussels 1000, Belgium.}

% Contact information of the corresponding author
\corres{Correspondence: polina.lemenkova@ulb.be; Tel.: +32471860459 (P.L.)}

% The commands \thirdnote{} till \eighthnote{} are available for further notes

%\simplesumm{} % Simple summary

% Abstract (Do not insert blank lines, i.e. \\) 
\abstract{This paper presents a cartographic framework based on algorithms of GMT codes for mapping seismically active areas in Venezuela. Data included raster grids from GEBCO, EGM-2008, and vector geological layers from the USGS. The data were processed iteratively in the console of GMT, converted by GDAL, formatted, and mapped for geophysical data visualisation; the QGIS was applied for geological mapping. We analysed 2,000 samples of the earthquake events obtained from the IRIS seismic database with 25-year time span (1997-2021) for mapping the seismicity. The approach to linking geological, topographic and geophysical data using GMT scripts aimed to map correlations between the geophysical phenomena, tectonic processes, geological setting, seismicity and earthquakes. Practical application of GMT scripts consists in automated mapping for visualization of geological risk and hazards in the mountainous region of the Venezuelan Andes. The proposed method integrates the approach of GMT scripts with state-of-the-art GIS techniques which demonstrated an effective tool for mapping spatial datasets and rapid data processing in iterative regime. In this context, using GMT and GIS for finding similarity between the regional earthquake distribution and geological and topographic setting is essential for hazard risk assessment. This study can serve as a basis for predictive seismic analysis in geological vulnerable regions of Venezuela. Besides technical demonstration of GMT algorithms, the study also contributes to the geologic \textcolor{blue}{ and }geophysical mapping\textcolor{blue}{,} and seismic hazard assessment in South America. We presented full scripts using for mapping in the GitHub repository.}

\keyword{Risk; hazard; sustainability; earthquake; seismicity; cartography; GMT; geophysics} 

\PACS{91.10.Da, 91.10.-v, 91.10.Jf, 91.10.Op, 91.30.Dk, 93.85.Pq, 91.70.-c}
\MSC{86Axx, 86A04, 86A15, 86A30, 86A60, 86-XX, 86-04, 86-08}
\JEL{Y91, Q00, Q01, Q2, Q20, Q24, Q3, Q35, Q5, Q50, Q51, Q54, Q55, Q56, C6, C61, C63}

%%%%%%%%%%%%%%%%%%%%%%%%%%%%%%%%%%%%%%%%%
\begin{document}

\section{Introduction}

\subsection{Background and motivation}

Nowadays, geologic hazards are acknowledged as one of the most disastrous and risk events for mountainous regions. These involve a series of negative consequences for nature (landslides, tsunami, rock falls, land mass movements) and complex social problems (destruction of houses and infrastructure, injures and losses of human lives) caused by seismic risks \citep{Laffaille,Schmitz2020,Weber}. Early prediction of seismic risks in geologically vulnerable regions has become more valuable in a large variety of \hl{the }time-critical geospatial applications focused on evaluating the geological hazards \hl{for mitigating of }earthquake hazards at the regional scale. The integration of geological and geophysical data is effective \hl{approach for }risk assessment in mountainous regions, preventing damages and dangers, ands to serve for early warning, long-term monitoring\hl{, management }and mitigation of seismic hazards. 

\hl{Recently, }advances in cartographic development and mapping made an important application of spatial data analysis become real: predicting and prognosis of seismic activities or imminent earthquake events from the observed events in the past \cite{KHAN2020101642,MESGAR2017454,NATH2005329}. \hl{With enough data on earthquake events, these methods }enable\hl{ an accurate }prognosis of seismicity aimed at mitigating hazard risks\hl{ }\cite{,KWAG20227230,KWAG2021107809}. Early detection of seismic risks can be solved by evaluating and visualising different geophysical determinants that affect disaster fatality at the country level, namely seismic hazard strength, population density\hl{ and }distribution of locations\hl{, }demographics and socioeconomics data, as well as geological background layers \cite{LIU20111084,ijerph191912351,rs14174269,ijerph191811469}. For instance, a systematic examination of the correlations and links among various geophysical and geological determinants underlying disasters \hl{is useful for }preventing\hl{ }socioeconomic risks in\hl{ }geologically unstable zones \cite{YOUSUF2020100064,KAMATA2022103253}. In the earthquake analysis, the capability of predicting \hl{earthquake} events, their magnitude and focal depth is highly desirable for real-time risk assessment \cite{SHAO2022100177,YAGHMAEISABEGH2022107173}\hl{.} 

To build\hl{ }effective geological\hl{ }modelling\hl{ approaches }for seismic hazard prediction,\hl{ }real-time mapping applications\hl{ are beneficial for earthquake} prognosis\hl{. }However, obtaining an {automated approach for }cartographic dataset \hl{processing }for seismic prediction\hl{ is usually difficult, resulting in biases to a particular GIS software. }The approach presented in this study aims to solve \hl{this }problem\hl{ }by introducing the integrative study on scripts developed in the Generic Mapping Tools (GMT) scripting toolset \cite{Wessel}\hl{ in addition to the traditional mapping}. Such a framework, which is different from the \hl{conventional state-of-the-art } GIS \hl{methods}, presents a \hl{rapid workflow of} mapping \hl{based on the data processing from the} console using\hl{ }command line and codes written by the GMT syntax. Although \hl{advanced methods of }geophysical mapping \hl{present }a very important task of such seismically active regions as South America, \hl{scripting cartography remains }a quite new topic. To the best of our knowledge, the majority of similar studies use the conventional GIS rather than scripts for \hl{geophysical }data visualisation \cite{Castillo}\hl{.}

While conventional methods of GIS enable to plot the maps in a traditional regime, this often requires manual operations and is not suitable for automated mapping. In such a way, it limits the types of cartographic activities that can be used for operative geological analysis and seismic prediction. Besides, data handling is limited by \hl{the }restricted compatibility of formats accepted in GIS\hl{. }In contrast, scripts present a rapid and repeatable procedure of automated mapping that can be integrated \hl{into }real-time risk assessment and hazard analysis. \hl{As a response to these needs, }GMT presents a more flexible \hl{cartographic }approach and\hl{ }versatile strategy of \hl{spatial }data processing, aimed at \hl{seismic }risk assessment\hl{, geologic exploration and }management using multi-source datasets. 

\subsection{Actuality and objectives}

The integrated mapping using various multi-source data has the objective to visualise a complex interplay among the geological, topographic, tectonic and geophysical\hl{ parameters. Regional effects of these setting make }the Caribbean region and\hl{ }Venezuela exposed to a high risk of earthquakes,\hl{ }volcanos and associated tsunami hazards. For instance, tectonic structure\hl{ of the }deep basins and erosional troughs affects submarine geomorphology\hl{ }which can be interpreted using geophysical and gravimetric measurements\hl{. }The earthquake-related aftershocks, hazards and\hl{ }landslides \hl{result in the damaged }geotechnical infrastructure and bridges \cite{Salcedo}, demolished buildings\hl{, }lost of human lives \cite{Rodriguez2018}, aseismic slips and creeps \cite{Pousse2016}\hl{. Such consequences }affect both nature and social \hl{sectors }of Venezuela and require regular monitoring \hl{and} mapping\hl{ for }risk and vulnerability assessment. 

\hl{The }visualized information of these processes is \hl{important }for seismic prognosis and earthquake modelling. Recent research demonstrated that geological and geophysical mapping creates a basis for natural hazard risk assessment \cite{su10020304,su131810232,Lemenkova202212,app11219919,ijgi10030118}. Therefore, the data derived from hazard risk mapping can be used as a background information and recognition of complex geological actions and processes\hl{. }Previous cartographic investigations have been carried out mostly using traditional GIS in geologic and environmental studies \citep{Gonzalez,Hackley,Schmitz2021,Sun}. However, none of them focused on the integrated approach \hl{which combines the }diverse cartographic techniques for geophysical mapping of Venezuela. 

\hl{Scripting technique applied to }geological mapping has not been adequately presented in the existing literature\hl{, compared to the traditional GIS}. At the same time, scripts\hl{ }employed for visualization of the complex geophysical and geologic phenomena enable automation\hl{ }and increase the precision of \hl{the }cartographic data processing. The use of \hl{the }machine-based methods enables to accurately map and highlight\hl{ }correlations between \hl{the }geological and geophysical phenomena\hl{. }With this aim, the GMT \hl{scripts }combined with QGIS and its plugins\hl{ }is applied for geophysical mapping of Venezuela, \hl{for visualization of }regional seismicity \hl{in connection }to the geologic structure \hl{and geophysical setting}. 

\hl{The }cartographic objective \hl{of this study }is to demonstrate the \hl{functionality }of GMT scripts in a combination with GIS techniques. The geological aim is to highlight the correlation between the geological setting, topography and geophysical anomalies and regional seismicity by the analysis of distribution, magnitude and depth of \hl{the }earthquakes. \hl{To this end, this }study applied the GMT scripting cartographic techniques and QGIS \hl{for processing }high-resolution open sources raster grids and vector data. We used the variety of datasets to build the project with observed geological and geophysical processes and structures\hl{. The primary }context \hl{of this work is }to analyse seismic events, \hl{from }the visualization of earthquake locations\hl{ }in the region. 

\section{Study Area}
This study focuses on Venezuela\hl{, }one the of the most seismically active region of \hl{South America and in the }world.  (Figure \ref{fig1}). 

\begin{figure}[H]\centering
	\includegraphics[width=11.0 cm]{Fig_01}
	\caption{Topographic map of the study area comprising the two countries of South America: the Venezuela and the Trinidad and Tobago. Mapping: GMT. Source: authors. \label{fig1}}
\end{figure}  

\hl{The }location in the Andes\hl{, makes Venezuela }exposed to the high risk of seismic hazards.
Geologically, the most important factors that have a significant influence on the seismicity in Venezuela and surroundings include lithosphere plate subduction and associated complex geological processes\hl{ of the }oblique collision of the Caribbean arc and a zone of active deformation in the eastern offshore Trinidad area \cite{ESCALONA20118}. This is reflected in the formation of the Venezuelan Andes as a N50°E-oriented mountain belt extending from the Colombian border in the SW to the Caribbean Sea in the NE with associated Boconó and Valera thrust faults \cite{DHONT2012245}. Seismotectonic activities in Venezuela can be characterised by a spatio-temporal composition of the constituent geophysical and geological forces\hl{, their }evolutional processes and\hl{ }complex interactions. Such linkages are related to the actual location and distribution of seismic sources of the earthquakes in the Andean region, with regard to the\hl{ }tectonics and overall dynamics of the Earth's crust in South America. 

The \hl{high }frequency and magnitude of earthquakes in Venezuela are associated with \hl{complex }geologic and geophysical setting\hl{, Figure} \ref{fig2}. \hl{Besides, }tectonic interaction of the lithosphere plates\hl{ triggers }seismic activity and tectonic-related hazards \hl{which }affect the lives of people and infrastructure \citep{Masy,Perez2018,Suarez}.\hl{ }The Caribbean-South American plate boundary is a very wide tectonically active zone \cite{Audemard1993,Backe} with \hl{recorded }seismicity at various depths. \hl{Thus}, anomalous earthquake series took place in the east coast of the Trinidad with reported depth of the main shock, 53 km \cite{Marshall}. Seismicity deeper than 40 km is reported in the southern end of the Lesser Antilles subduction zone and\hl{ }active island arc \cite{Reinoza}. 

\begin{figure}[H]\centering
	\includegraphics[width=14.0 cm]{Fig_02}
	\caption{Geologic map of Venezuela, Trinidad and Tobago and neighboring countries. Mapping: QGIS. Source: authors.\label{fig2}}
\end{figure}  

\hl{Venezuela} is one of the most vulnerable countries to earthquakes\hl{ in the Amazon Basin}. \hl{The }geological hazards \hl{in this region }are exacerbated by social vulnerability, due to the anthropogenic exposure with\hl{ }frequently injured and affected population \hl{living }in \hl{the }high mountain regions of Andes\hl{, which }increases the population at risk during earthquakes\hl{. }Furthermore, catastrophic earthquakes produce a chain effect triggering tsunami in succession, which amplifies the damage to the environment and human society. For example, it is reported that since the 19th century\hl{, }the tsunami in the Caribbean Sea resulted in several thousands of lives lost and many people exposed at risk and damage from the landslides \cite{Harbitz}. 

Venezuela is\hl{ }located in the north of South American continent where a collision of South American and Caribbean tectonic plates\hl{ causes }complex regional geologic patterns\hl{ }(Figure \ref{fig2}). The geodynamics of the interacted plates results in the formation of \hl{the }thrust belts and foreland basins in northern Venezuela that are \hl{shaped }by the interplay of geologic layers.\hl{ }The outcrops of \hl{the }Precambrian (pC), Quaternary (Q), Triassic (T), Paleozoic-Mesozoic intrusives (PMi) are dominating in the central regions of the country\hl{, }while regional outcrops of \hl{the }Cretaceous (K), Paleozoic metamorphics (PZm) and Mesozoic-Cenozoic intrusives (MCi) are \hl{found }in the Falcón Basin (Figure \ref{fig2}). 

The north-western and north-eastern regions\hl{ of Venezuela are }prone to earthquakes \cite{Audemard,Perez} due to the tectonic activity of the lithosphere plates \citep{Barr,Hayes}. \hl{According to IRIS database, the }repetitive occurrence of seismic events, such as volcanism and earthquakes in north-western\hl{ }Venezuela \hl{as well as} Trinidad and Tobago are recorded as \hl{a }series of\hl{ }events \hl{in }recent 25 years. The ongoing effects of earthquakes in Venezuela \citep{Audemard2006,Baumbach,Colon,Pousse2018,Rial,Yamazaki} and coastal regions of Trinidad and Tobago \citep{Deville,Latchman}, made geophysical mapping important task \hl{for} seismic hazard mapping, earthquake \hl{prognosis }and risk assessment \cite{Bucci}.

\begin{figure}[H]\centering
	\includegraphics[width=14.0 cm]{Fig_03}
	\caption{Geologic provinces in Venezuela, Trinidad and Tobago and neighbouring countries. Mapping: QGIS. Source: authors.\label{fig3}}
\end{figure} 

The Caribbean Sea region is notable for complex geologic setting (Figure \ref{fig3}) where several tectonic plates experience collision, interaction and subduction: Caribbean, Cocos, Nazca and South America plates. This results in high seismic activity in the region under the influence of tectonic processes which his reflected in current topography, geology and geomorphology of South America in general and Venezuela in particular, as has been discussed in previously published relevant papers \citep{Diebold,Fajardo,Lemenkova2020b,Perez1981}. 

The geodynamic heterogeneity of the upper mantle resulted in specific mountain geomorphology with notable features such as the extent of the Tobago Trough, Lara-Falcón Basin and crustal structure of the Cordillera de Mérida in southwestern Venezuela, as important segments of Andean orogeny \cite{AVILAGARCIA2022103853}. 
The distribution of earthquakes in this region is strongly controlled by the tectonic and geophysical setting, i.e., closeness to the margins of the colliding South America and Caribbean tectonic plates. Another key factor is a mountain geomorphology and topographic variability shaped by the geophysical-geological forces which affect the variations in depth and magnitude of seismic events. 

Tectonic active movements cause the formation of the orogenic mountain belt and topography, bathymetry complicated with deep-sea trenches and troughs in the margins of the Pacific Ocean and the Caribbean Sea along the subduction path \citep{Gough,Lemenkova2020a}. Complex geological history is especially notable in the NW Venezuela, Falcón Basin, where the interactions between the lithosphere plates and subducting Caribbean slab cause the associated seismic activity and geologic processes: outcrop of igneous intrusive bodies, crustal thinning, decreased Moho depth, gravity anomalies, rifting \cite{Bezada}. \cite{Mazuera} identified crustal thinning beneath the Falcón Basin along the western extension of the Oca-Ancón fault system which they interpreted as a back-arc basin, and suture zones between the Proterozoic and Paleozoic provinces (Ouachita-Marathon-related suture), and Paleozoic and Meso-Cenozoic terranes (peri-Caribbean suture) where seismic velocity changed. 

%%%%%%%%%%%%%%%%%%%%%%%%%%%%%%%%%%%%%%%%%%
\section{Materials and Methods}

The study area encompasses the spatial extent of 59,5°-74°W, 0°-12.5°N (Figure \ref{fig1}). The topographic, geophysical, seismic and geologic data sources were used to present a multi-disciplinary cartographic analysis of Venezuela. The topographic data were collected through the General Bathymetric Chart of the Oceans (GEBCO) open access repository with NetCDF raster grid \cite{GEBCO} which presents the most comprehensive data for topography and bathymetry of the Earth \cite{Schenke}. The GEBCO data grid is widely used in geosciences due to its unprecedented resolution and precision \citep{Mcmichael,Cortina}. 

The geological data (Figures \ref{fig2} and \ref{fig3}) have been plotted using the QGIS, because this open source software is compatible with the ArcGIS data format that have been used in geologic maps as a .shp vector layers. The visualization of these maps followed the traditional mapping process of the QGIS \cite{Lemenkova2022a}. In the first stage, vector layers of the geologic provinces and outcrops were visualized with in a Graphical User Interface (GUI). In the second stage, the study area was designated using the overlay of the Digital Chart of the World (DCW), which is a comprehensive 1:1,000,000 scale vector base map uploaded into the QGIS project. 

Maps of topographic, seismic and geophysical data (Figures \ref{fig1}, \ref{fig4}, \ref{fig5}, \ref{fig6} and \ref{fig7} were made using scripting technique of the GMT \cite{Wessel} following the existing examples of scripting \citep{Lemenkova2022b,Lemenkova202217}. The methodological workflow included 5 algorithms by GMT codes \hl{is presented in the Appendix section of this article.} Prior to data visualization and mapping, the GMT scripts were written using its syntax and organized in the Xcode environment. The cocding language of GMT was used during scripting and several modules applied for plotting various cartographic elements. the area of interest was first selected and cut from the GEBCO grid using the following code: $'gmt grdcut GEBCO\_2019.nc -R286/300.5/0/12.5 -Gve\_relief.nc'.$

Following the first step, all other the substantial cartographic elements (image visualized in selected color palette, grid, clipped insert region in the globe map, coastlines, borders, rivers, barlegend, scale bar, text annotations and time stamp) were added in a script using the following GMT modules: psbasemap, grdimage, psscale, pstext, psclip, grdcontour. The technique of scripting was applied from the existing detailed descriptions \cite{Lemenkova2021b}. For example, the 'grdcontour' module was used for computing, modelling and visualizing topographic isolines using quantitative DEM data as follows: $'gmt grdcontour ve\_relief1.nc -R -J -C1000 -Wthinner,darkbrown -O -K >> \$ps'$. The purpose of the geoid visualization (Figure \ref{fig4}) in relation to geophysical differentiation was to investigate how gravity values respond to the effects of topographic elevation and bathymetric depressions. Data on geoid undulations have been collected from the EGM2008 \cite{Pavlis}. 

Whilst GEBCO data were gathered in NetCDF format ($GEBCO\_2019.nc$), the geoid data needed to be converted first from the .adf format which was done by the following code of GMT: $'gmt grdconvert n00w90/w001001.adf geoid\_02.grd'.$ The data extent and range were examined by GDAL and interpreted for the color palette as follows: $'gmt makecpt -C33_blue\_red.cpt -T-55/25 > colors.cpt'.$ The geoid was modelled using the following code: $'gmt grdimage geoid\_02.grd -Ccolors.cpt -R286/300.5/0/12.5 -JM6.5i -P -Xc -I+a15+ne0.75 -K > \$ps'.$

Geophysical data of gravity grids (Figure \ref{fig5} and \ref{fig6}) were acquired from available sources \cite{Sandwell}. The seismic data were downloaded from the IRIS database using spatial extent of the Venezuela, Trinidad and Tobago and surroundings, respectively (Figure \ref{fig7}). The gravity grids were also converted from the initial raw IMG format to the GRD format compatoble with the GMT as follows: $'gmt img2grd grav\_27.1.img -R286/300.5/0/12.5 -Ggrav\_VE.grd -T1 -I1 -E -S0.1 -V'$ (here for Fig. 5). The same procedure has been repeated for the Fig. 6 (vertical gradient of gravity) and followed the example of published work \cite{Lemenkova2021d}. The annotations have been added on all the maps using the 'pstext' module of GMT as in the following example: $'gmt pstext -R -J -N -O -K -F+f10p,0,black+jLB+a-0 >> \$ps << EOF 294.20 8.24 Ciudad Bolívar EOF'.$ The data were converted to ‘ngdc’ format, visualized by 'psxy' module and assessed for seismicity. 

The earthquakes on the Figure \ref{fig7} (map of seismicity) have been plotted using the 'psxy' module of GMT by the following code: $'gmt psxy -R -J quakes\_VE.ngdc -Wfaint -i4,3,6,6s0.05 -h3 -Scc -Csteps.cpt -O -K >> \$ps'.$ Here the ‘i4,3,6,6s0.05’ means the numbers of column in the table ($quakes\_VE.ngdc$) that was converted from the original CSV from the IRIS. The volcanoes were visualized by the following code: $‘gmt psxy -R -J volcanoes.gmt -St0.4c -Gred -Wthinnest -O -K >> \$ps’$. The complex legend showing the magnitude of events was plotted using the following code: $‘gmt pslegend -R -J -Dx1.5/-3.0+w17.8c+o-2.0/0.1c  -F+pthin+ithinner+gwhite -O -K << FIN >> \$ps H 10 Helvetica Seismicity: earthquakes magnitude (M) from \\1.9 to 7.3 N 9 S 0.3c c 0.3c red 0.01c 0.5c M (7.1-7.3) <...> FIN’.$

The full algorithms of GMT codes are available in author’s GitHub repository: \\\href{https://github.com/paulinelemenkova/Mapping_Venezuela_GMT_Scripts}{https://github.com/paulinelemenkova/Mapping\_Venezuela\_GMT\_Scripts}. 

%%%%%%%%%%%%%%%%%%%%%%%%%%%%%%%%%%%%%%%%%%
\section{Results}

Figure \ref{fig1} has been plotted using GEBCO raster grid selected due to its quality and reliability. Precise topography is fundamental to the understanding of geophysics and tectonics and associated effects on ocean and terrestrial geomorphology, mapping seismicity, earthquake localization, risk assessment and tsunami wave propagation. Here the topographic and bathymetric elevations of the study area encompassing Venezuela, Trinidad and Tobago islands range from -4948 m to 5536 m according to the GEBCO grid. 

Figure \ref{fig2} and \ref{fig3} indicate spatial extent of geological provinces and outcrops in Venezuela. The map in Figure \ref{fig2} shows that most of the territory of Venezuela is covered by Guyana Shield (aquamarine color) in the south-eastern parts of the country. The Barinas-Apure sedimentary basin is located in the Llanos region and the Andean foothills region in the western Venezuela \cite{Callejon}. The Barinas-Apure basin (yellow color in the map) includes oil fields that is the 3rd most important producing area in Venezuela after the Maracaibo and Eastern Venezuela basins according to the petroleum accumulation. The Falcón sedimentary basin, which produces indigenous Miocene oil from structural and stratigraphic accumulations \cite{Rodriguez} is coloured slate blue in the map. 

The South Caribbean deformed belt (purple colour in Figure \ref{fig2}) represents a submarine prism formed between the subducting tectonic plates in the Colombian and Venezuelan basins and arc terranes along the northern rim of South America \cite{Mann}. Its SW part, the Sinu fold belt, forms an accretionary wedge in the place of the subduction of the Caribbean Plate under South American Plate \cite{Rodriguez2021}. Other important geological provinces include the Guyana Basin which stretches along the passive margin of NE South America and the Maracaibo Basin (green, Figure \ref{fig2}). The Cariaco Basin (beige color in Figure \ref{fig2}) is an East-West stretching basin \cite{Schubert} situated in the Gulf of Cariaco, on the continental shelf of the Caribbean Sea, off the eastern coast (near the Barcelona, see Figure \ref{fig1}). 

The Tobago Trough (chartreuse green colour in Figure \ref{fig2}) is a modern marine forearc basin of the Lesser Antilles arc located above the sedimentary succession at least 10 km thick and probable oceanic basement and formed by sediments from various stratigraphic sequences: Quaternary, early Pleistocene – late Miocene, Miocene and late Eocene and may also include mid-Cretaceous or older \cite{Speed}. The Quaternary sediments covering the most of the north-west of the country (dark pink colour in Figure \ref{fig3}) corresponds with the previous studies on Plio-Quaternary extension in the Venezuelan Andes showing the extensional structures corresponding to the elongated tilted blocks with geometry and kinematics of the structures corresponding to the earlier syn-orogenic extension \cite{Dhont}.

\begin{figure}[H]\centering
	\includegraphics[width=12.0 cm]{Fig_04}
	\caption{Geoid model over the Venezuela, Trinidad and Tobago. Mapping: GMT. Source: authors.\label{fig4}}
\end{figure}  

Figure \ref{fig4} show the variations of the geoid heights over the study area. The highest values of geoid (up to 25 m, bright red colors in Figure \ref{fig4}) are visible in the south-western region of the map (Colombia) while the majority of Venezuela covers the extent of -50 to 8 m with the highest values in the mountains and clearly visible correlation with topographic elevations. Thus, if compared geoid mode in Figure \ref{fig4} with topographic and geologic maps in Figure \ref{fig1} and \ref{fig2}, respectively, the elevated values of geoid are notable in the region of the Guyana Shield. The Guyana Shield is a Precambrian geological formation in the NE coast of the South America and an area of special importance in historical geology of Venezuela due to high potential for mineral exploration \cite{Gibbs}. 

The environmental importance of this unique region consists in the highest biodiversity in the world including many endemic species \cite{Hollowell} in tropical forests of Amazonia \cite{Hammond}. The undulations of geoid in the Guyana Highlands corresponds to the higher topographic elevations on the shield as table-like mountains ‘tepuis’. The marine areas of the Caribbean Sea have in general lower geoid heights (blue colored areas in Figure \ref{fig4}) which well corresponds to the bathymetric depressions reflected in the geophysical setting of the Earth.

Figure \ref{fig5} shows variations in the gravity over the study area. The notable free-air gravity low coincides with distribution of the Guyana basin, SE part of the Venezuela, Trinidad and Tobago, which points at the tectonic basement depression formed during Jurassic rifting, and Lake Maracaibo, an important geologic region of Venezuela. The NE coasts of the Lake Maracaibo are occupied by the Bolivar Coastal Fields, the largest oil field in South America, the 2nd oil field most prolific hydrocarbon basins in the world after the Middle East with basin area of 50,000 km2 and oil production of 30 *109 bbl \cite{Escalona}. 

\begin{figure}[H]\centering
	\includegraphics[width=12.0 cm]{Fig_05}
	\caption{Free-air gravity model of the Venezuela, Trinidad and Tobago. Mapping: GMT. Source: authors.\label{fig5}}
\end{figure}   

The Maracaibo Lake has a wide variety of oil reservoir rocks with structural traps in the tectonic units, e.g. normal or inverted faults on the South American plate \cite{Bockmeulen}. The oil formed in the Maracaibo basin is mostly heavy black oil with asphaltene deposition \cite{Escobar,Memon}. The gravity in the Lake Maracaibo reach below -200 mGal correlating to the bathymetric depression (Figure \ref{fig1}) and geologic extend of the Maracaibo Basin (Figure \ref{fig2}).\hl{ The }highest values for gravity reach +200 mGal and well corresponds to the topographic elevation of the Cordillera de Mérida, Columbian Andes, Guiana Highlands and Sierra Parima. The Amazon and Orinoco Flow basin mostly have gentle fluctuations of gravity with values in the range between -25 to 30 mGal (Figure \ref{fig5}).

Figure \ref{fig6} refers to a vertical gravity gradient of normal gravity as an approximation in which a gravity measurement point, which is almost never placed on the reference ellipsoid’s surface is adjusted by geophysical corrections \citep{Li,Yilmaz}. The details regarding the distribution of vertical gravity values over the Venezuela and the Trinidad and Tobago showed that the lowest values (below -100 mGal) mostly correspond to the extreme amplitude of topographic changes (e.g. abrupt border between the mountain hills and margins of the bathymetric depression) whilst the highest values (over 100 mGal) were detected in the top of the mountains, well corresponding to the local geomorphology of the Venezuela and the Trinidad and Tobago region.

\begin{figure}[H]\centering
	\includegraphics[width=12.0 cm]{Fig_06}
	\caption{Vertical gradient of gravity in the Venezuela, Trinidad and Tobago. Mapping: GMT. Source: authors.\label{fig6}}
\end{figure}  

The distribution of the earthquakes and seismicity of the Venezuela and the Trinidad and Tobago (Figure \ref{fig7}) depends on the proximity of area to the lithosphere plate border with active margin such as South American and Caribbean plates (depicted in the map as purple thick line). The second factor includes the geologic active processes (faults), and geomorphology. Hence, Cordillera de Mérida has the majority of the detected earthquakes following the Trinidad and Tobago active zone. The majority of mapped earthquakes\hl{ in }Venezuela \hl{belongs} to the category ‘shallow’, that is, having depth between 0 and 70 km. Earthquakes at depth 74 to 155 km are located in the offshore Sucre, Venezuela, those detected in the Trinidad and Tobago are dominantly shallow not exceeding 70 km. The deepest earthquakes from the inspected database are located in the northern Columbia (182-200 km).

\begin{figure}[H]\centering
	\includegraphics[width=12.0 cm]{Fig_07}
	\caption{Seismic map of earthquakes in the region of Venezuela, Trinidad and Tobago. Mapping: GMT. Source: authors.\label{fig7}}
\end{figure}   

\section{Conclusion}

In this paper, we have proposed two algorithms of geophysical mapping of Venezuela which include GMT based scripts and traditional plotting using QGIS aimed at geohazard risk analysis of Venezuela. Geohazard investigation and risk assessment require advanced methods of mapping which are based on programming algorithms, as demonstrated in this paper. The operative use of the correctly and effectively plotted maps, along with tabular data, supports reasonable and argued recommendations aimed at prevention and mitigation of earthquake hazard risks and seismically active region, such as South American Andes. The algorithms of GMT enabled improved performance of cartographic workflow than conventional approaches using the software. 

We demonstrated technical advantages of programming approaches in cartography, which include cartographic scripting, a rapidly developing branch of geo-information. Due to the high degree of automation, scripting is the best technical cartographic solution for seismic mapping, because it enables rapid machine-based processing of data, which is useful for prognosis, predictive mapping, risk assessment and hazard visualization in real time regime. Automation enables to process big datasets rapidly and effectively using multi-source heterogeneous data. Besides, that, known advantages of scripting cartography include machine-based graphics that results in an aesthetically refined plotting by GMT. 

Scripting cartographic methods of GMT were addressed to evaluate the geophysical and geologic correspondence in a seismically active region of Venezuela and the Caribbean Sea basin. We used high-resolution geophysical, topographic and geological datasets, such as GEBCO, EGM-2008 for geoid, gravity grids and seismic data obtained from IRIS. The results were corresponding and supporting previous studies on geological and geophysical methods of gravity measurements, seismic observations and earthquake monitoring in Venezuela and the surroundings \citep{Guenther,Schmitz2008}.  

The study presented the machine-based mapping workflow which contributed to the regional studies of the eastern Caribbean and Venezuela. The GMT code snippets are presented in the GitHub repository, as a demonstration of the technical possibilities of this toolset. The results showed optimised cartographic performance and demonstrated the advantages of the automating processing of geophysical datasets. The multi-source geophysical, geologic, seismic and topographic data were used for comparative analysis of the geophysical processes. The tectonic evolution of the Caribbean Sea basin largely affected its geological setting that can be tracked in relevant geophysical data which indicate high seismicity in the Andes and high earthquake risk.

Since eastern Caribbean is a tectonically complex area which includes a subduction zone, island arc and active volcanoes \cite{Ambeh}, regional analysis and visualization of the varied setting was performed. Specifically, high seismicity was detected in the north-west of Venezuela, and north-east region of Trinidad and Tobago, which is consistent with geologic structure associated with tectonic setting, and related movements of the Caribbean Plate subduction. Although continued research is needed to improve earthquake monitoring in Venezuela region, this study presented a contribution to the seismic studies in the Caribbean and South America, showing earthquakes detected according to the IRIS database based on the 25-year time span (1997-2021). 

A better understanding of geological and tectonic factors affecting seismicity and regulating the appearance, magnitude and focal depth of the earthquakes is critical to seismically active regions, such as the Caribbean and Venezuela. Practical applications of geologic risk assessment in South America include the use of information by public administration, city management and urban planning for environmental management and hazard risk mitigation and management in regions prone to geological hazards, such as Venezuela. 

\hl{The main contributions we have made and the actuality of the presented study are highlighted as follows:}

\begin{enumerate}
  \item \textbf{Region} The region of Venezuela is at high risk of seismicity and earthquakes, \hl{because it is }located in the zone of the\hl{ collision of the }South America and Caribbean tectonic plates and Andean orogeny. This requires spatial analysis of the regions at risk.
  \item \textbf{Geohazards} Risk assessment of geophysical hazards is based on\hl{ the effective }integration and visualisation of the multi-source data\hl{, }which provide detailed insights into the environmental and geological settings of Venezuela and support complex investigations in a seismically active regions of South America. 
    \item \textbf{Data} Geological disasters are related to high seismicity, exposure to hazards, and vulnerability of people at risk. This requires complex detailed studies summarising these factors and visualising regional geophysical setting\hl{, as shown in this paper}.
   \item \textbf{Methods} Combination of scripting GMT and GIS is \hl{a solid foundation} for cartographic data analysis. \hl{The efficient method of }multi-format data processing ensures hazard mapping, geophysical and geological visualization, \hl{aimed at} preventing and evaluating the seismic hazards.
\end{enumerate}

%%%%%%%%%%%%%%%%%%%%%%%%%%%%%%%%%%%%%%%%%%
\vspace{6pt} 

\authorcontributions{Supervision, conceptualization, methodology, software, resources, funding acquisition, project administration, O.D; writing---original draft preparation, methodology, software, data curation, visualization, formal analysis, validation, writing---review and editing, investigation, P.L. All authors have read and agreed to the published version of the manuscript.}

\funding{\textcolor{blue}{This project was supported by The Federal Public Planning Service Science Policy or Belgian Science Policy Office, Federal Science Policy - BELSPO (B2/202/P2/SEISMOSTORM)}.}

\institutionalreview{Not applicable}

\informedconsent{Not applicable}

\dataavailability{The GitHub repository contains GMT scripts using for mapping in this study: \href{https://github.com/paulinelemenkova/Mapping_Venezuela_GMT_Scripts}{https://github.com/paulinelemenkova/Mapping\_Venezuela\_GMT\_Scripts}} 

\acknowledgments{The authors thank the three anonymous reviewers for careful reading, critical suggestions and useful comments that helped improve an earlier version of this manuscript.}

\conflictsofinterest{The authors declare no conflict of interest.} 

\abbreviations{Abbreviations}{
The following abbreviations are used in this manuscript:\\

\noindent 
\begin{tabular}{@{}ll}
CSV & comma-separated values \\
DCW & Digital Chart of the World \\
EGM & Earth Gravitational Models \\ 
GEBCO & The General Bathymetric Chart of the Oceans \\
GIS & Geographic Information System \\
GMT & Generic Mapping Tools \\
GUI & Graphical User Interface \\
QGIS & Quantum GIS \\
IRIS & Incorporated Research Institutions for Seismology \\
NetCDF & Network Common Data Form \\
NGDC & National Geophysical Data Center \\
USGS & United States Geological Survey \\
\end{tabular}}

%%%%%%%%%%%%%%%%%%%%%%%%%%%%%%%%%%%%%%%%%%
%% Optional
\appendixtitles{yes} % Leave argument "no" if all appendix headings stay EMPTY (then no dot is printed after "Appendix A"). If the appendix sections contain a heading then change the argument to "yes".
\appendixstart
\appendix
\section[\appendixname~\thesection]{GMT scripts}
\subsection[\appendixname~\thesubsection]{GMT script for topographic mapping}

\begin{lstlisting}[language=bash, caption=\textbf{GMT code used for plotting Figure \ref{fig1} (topographic map)}]
#!/bin/sh
# Purpose: shaded relief grid raster map from the GEBCO dataset (here: Venezuela)
# GMT modules: gmtset, gmtdefaults, grdcut, makecpt, grdimage, psscale, grdcontour, psbasemap, gmtlogo, psconvert
# http://soliton.vm.bytemark.co.uk/pub/cpt-city/arendal/tn/arctic.png.index.html
# GMT set up
gmt set FORMAT_GEO_MAP=dddF \
# Overwrite defaults of GMT
gmtdefaults -D > .gmtdefaults
#chsh -s /bin/bash
chsh -s /bin/zsh
gmt grdcut GEBCO_2019.nc -R286/300.5/0/12.5 -Gve_relief.nc
gmt grdcut ETOPO1_Ice_g_gmt4.grd -R286/300.5/0/12.5 -Gve_relief1.nc
gdalinfo ve_relief.nc -stats
# Minimum=-4947.963, Maximum=5535.625, Mean=162.220, StdDev=788.097
# create mask of vector layer from the DCW of country's polygon
gmt pscoast -R286/300.5/0/12.5 -Dh -M -EVE > ve.txt
# Make color palette
gmt makecpt -Carctic.cpt > pauline.cpt
# Generate a file
ps=Topography_VE.ps
# Make background transparent image
gmt grdimage ve_relief.nc -Cpauline.cpt -R286/300.5/0/12.5 -JM6i -P -I+a15+ne0.75 -t20 -Xc -K > $ps
# Add isolines
gmt grdcontour ve_relief1.nc -R -J -C500 -W0.1p -O -K >> $ps
# Add coastlines, borders, rivers
gmt pscoast -R -J -P \
    -Ia/thinner,blue -Na -N1/thickest,darkred -W0.1p -Df -O -K >> $ps
# CLIPPING
# 1. Start: clip the map by mask to only include country
gmt psclip -R286/300.5/0/12.5 -JM6.0i ve.txt -O -K >> $ps
# 2. create map within mask
# Add raster image
gmt grdimage ve_relief.nc -Cpauline.cpt -R286/300.5/0/12.5 -JM6.0i -I+a15+ne0.75 -Xc -P -O -K >> $ps
# Add isolines
gmt grdcontour ve_relief1.nc -R -J -C1000 -Wthinner,darkbrown -O -K >> $ps
# Add coastlines, borders, rivers
gmt pscoast -R -J \
    -Ia/thinner,blue -Na -N1/thicker,tomato -W0.1p -Df -O -K >> $ps
# 3: Undo the clipping
gmt psclip -C -O -K >> $ps
# Add color barlegend
gmt psscale -Dg286/-1.0+w15.2c/0.4c+h+o0.0/0i+ml -R -J -Cpauline.cpt \
    -Bg500a1000f100+l"Color scale 'arctic': global bathymetry/topography relief [-5000 to 4000, mixed, RGB, 110 segments]" \
    -I0.2 -By+lm -O -K >> $ps    
# Add grid
gmt psbasemap -R -J \
    --MAP_FRAME_AXES=WEsN \
    --FORMAT_GEO_MAP=ddd:mm:ssF \
    -Bpx4f2a2 -Bpyg4f2a2 -Bsxg4 -Bsyg2 \
    -B+t"Topographic map of Venezuela, Trinidad and Tobago" -O -K >> $ps    
# Add scalebar, directional rose
gmt psbasemap -R -J \
    -Lx12.7c/-2.2c+c50+w200k+l"Mercator projection. Scale: km"+f \
    -UBL/-5p/-65p -O -K >> $ps
# Texts
gmt pstext -R -J -N -O -K \
-F+jTL+f10p,21,darkred+jLB+a-0 -Gwhite@60 >> $ps << EOF
290.1 11.8 Peninsula
EOF
# insert map
# Countries codes: ISO 3166-1 alpha-2. Continent codes AF (Africa), AN (Antarctica), AS (Asia), EU (Europe), OC (Oceania), NA (North America), or SA (South America). -EEU+ggrey
gmt psbasemap -R -J -O -K -DjBL+w3.3c+stmp >> $ps
read x0 y0 w h < tmp
gmt pscoast --MAP_GRID_PEN_PRIMARY=thinner,white -Rg -JG293/6N/$w -Da -Gbrown -A5000 -Bg -Wfaint -ESA+gpeachpuff -EVE+gyellow -Slightsteelblue3 -O -K -X$x0 -Y$y0 >> $ps
#gmt pscoast -Rg -JG12/5N/$w -Da -Gbrown -A5000 -Bg -Wfaint -ECM+gbisque -O -K -X$x0 -Y$y0 >> $ps
gmt psxy -R -J -O -K -T  -X-${x0} -Y-${y0} >> $ps
# Add GMT logo
gmt logo -Dx6.2/-2.9+o0.1i/0.1i+w2c -O -K >> $ps
# Add subtitle
gmt pstext -R0/10/0/15 -JX10/10 -X0.5c -Y5.0c -N -O \
    -F+f10p,Helvetica,black+jLB >> $ps << EOF
2.3 13.3 SRTM/GEBCO 15 arc sec resolution global terrain model grid
EOF
# Convert to image file using GhostScript
gmt psconvert Topography_VE.ps -A0.5c -E720 -Tj -Z
\end{lstlisting}

\subsection[\appendixname~\thesubsection]{GMT script for geoid mapping}
\begin{lstlisting}[language=bash, caption=\textbf{GMT code used for plotting Figure \ref{fig4} (geoid modelling)}]
#!/bin/sh
# Purpose: geoid of Venezuela
# GMT modules: gmtset, gmtdefaults, grdcut, makecpt, grdimage, psscale, grdcontour, psbasemap, gmtlogo, psconvert
# http://soliton.vm.bytemark.co.uk/pub/cpt-city/kst/tn/33_blue_red.png.index.html
# GMT set up
gmt set FORMAT_GEO_MAP=dddF \
# Overwrite defaults of GMT
gmtdefaults -D > .gmtdefaults
gmt grdconvert n00w90/w001001.adf geoid_02.grd
gdalinfo geoid_02.grd -stats
# Minimum=-70.856, Maximum=32.958, Mean=-24.768, StdDev=19.081
# Generate a color palette table from grid
gmt makecpt -C33_blue_red.cpt -T-55/25 > colors.cpt
# Generate a file
ps=Geoid_VE.ps
gmt grdimage geoid_02.grd -Ccolors.cpt -R286/300.5/0/12.5 -JM6.5i -P -Xc -I+a15+ne0.75 -K > $ps
# Add shorelines
gmt grdcontour geoid_02.grd -R -J -C1.0 -A2.0+f8p,0,black -Wthinner,dimgray -O -K >> $ps
# Add grid
gmt psbasemap -R -J \
    --MAP_FRAME_AXES=WEsN \
    --FORMAT_GEO_MAP=ddd:mm:ssF \
    -Bpx4f2a2 -Bpyg4f2a2 -Bsxg4 -Bsyg2 \
    --MAP_TITLE_OFFSET=0.8c \
    --FONT_ANNOT_PRIMARY=7p,0,black \
    --FONT_LABEL=7p,25,black \
    --FONT_TITLE=13p,25,black \
    -B+t"Geoid model (EGM-2008) of Venezuela and Trinidad and Tobago" -O -K >> $ps
# Add legend
gmt psscale -Dg286/-1.0+w16.5c/0.4c+h+o0.0/0i+ml+e -R -J -Ccolors.cpt \
    -Bg2f0.2a4+l"Color scale '33_blue_red': colour tables of IDL from the KDE scientific plotting tool [256, discrete, RGB, 252 segments, -T-55/25]" \
    -I0.2 -By+lm -O -K >> $ps
# Add scale, directional rose
gmt psbasemap -R -J \
    --FONT=7p,0,black \
    --FONT_ANNOT_PRIMARY=6p,0,black \
    --MAP_TITLE_OFFSET=0.1c \
    --MAP_ANNOT_OFFSET=0.1c \
    -Lx14.7c/-2.3c+c50+w200k+l"Mercator projection. Scale (km)"+f \
    -UBL/-5p/-65p -O -K >> $ps
# Add coastlines, borders, rivers
gmt pscoast -R -J -P -Ia/thinnest,blue -Na -N1/thick,brown -Wthick,darkslategray -Df -O -K >> $ps
# Texts
gmt pstext -R -J -N -O -K \
-F+jTL+f10p,21,darkred+jLB+a-0 >> $ps << EOF
290.0 10.5 (Coro)
EOF
# Add GMT logo
gmt logo -Dx7.0/-2.9+o0.1i/0.1i+w2c -O -K >> $ps
# Add subtitle
gmt pstext -R0/10/0/15 -JX10/10 -X0.1c -Y5.9c -N -O \
    -F+f10p,25,black+jLB >> $ps << EOF
3.0 13.6 World geoid image EGM2008 vertical datum 2.5 min resolution
EOF
# Convert to image file using GhostScript
gmt psconvert Geoid_VE.ps -A1.0c -E720 -Tj -Z
\end{lstlisting}

\subsection[\appendixname~\thesubsection]{GMT script for free-air gravity mapping}
\begin{lstlisting}[language=bash, caption=\textbf{GMT code used for plotting Fig. \ref{fig5} (free-air gravity modelling)}]
#!/bin/sh
# Purpose: mapping free-air gravity anomaly of Venezuela
# GMT modules: gmtset, gmtdefaults, img2grd, makecpt, grdimage, psscale, grdcontour, psbasemap, gmtlogo, psconvert, pscoast
# http://soliton.vm.bytemark.co.uk/pub/cpt-city/pj/4/index.html
# GMT set up
gmt set FORMAT_GEO_MAP=dddF \
    MAP_FRAME_PEN=dimgray \
# Overwrite defaults of GMT
gmtdefaults -D > .gmtdefaults
gmt img2grd grav_27.1.img -R286/300.5/0/12.5 -Ggrav_VE.grd -T1 -I1 -E -S0.1 -V
gdalinfo grav_VE.grd -stats
# Minimum=-218.668, Maximum=942.457, Mean=5.694, StdDev=64.118
# Generate a color palette table from grid
gmt makecpt -Chaxby -T-200/200 > colors.cpt
# Generate a file
ps=Grav_VE.ps
gmt grdimage grav_VE.grd -Ccolors.cpt -R286/300.5/0/12.5 -JM6.0i -P -I+a15+ne0.75 -Xc -K > $ps
# Add isolines
gmt grdcontour grav_VE.grd -R -J -C50 -A100 -Wthinner -O -K >> $ps
# Add grid
gmt psbasemap -R -J \
    --MAP_FRAME_AXES=WEsN \
    --FORMAT_GEO_MAP=ddd:mm:ssF \
    -Bpx4f2a2 -Bpyg4f2a2 -Bsxg4 -Bsyg2 \
    -B+t"Free-air gravity anomaly for Venezuela, Trinidad and Tobago" -O -K >> $ps 
# Add legend
gmt psscale -Dg286/-1.0+w15.0c/0.4c+h+o0.0/0i+ml+e -R -J -Ccolors.cpt \
    --FONT_LABEL=7p,Helvetica,black \
    --FONT_ANNOT_PRIMARY=7p,0,black \
    --FONT_TITLE=8p,25,black \
    -Bg50f5a50+l"Color scale 'haxby' B. Haxby's color scheme for geoid & gravity [C=RGB -T-200/200]" \
    -I0.2 -By+l"mGal" -O -K >> $ps
# Add scale, directional rose
gmt psbasemap -R -J \
    -Lx14.0c/-2.3c+c50+w200k+l"Mercator projection. Scale (km)"+f \
    -UBL/-5p/-70p -O -K >> $ps
# Add coastlines, borders, rivers
gmt pscoast -R -J -P -Ia/thinnest,blue -Na -N1/thick,white -Wthin,darkslategray -Df -O -K >> $ps
# Texts
gmt pstext -R -J -N -O -K \
-F+jTL+f14p,32,darkgreen+jLB+a-330 >> $ps << EOF
290.8 7.8 Los Llanos
EOF
# Add GMT logo
gmt logo -Dx7.0/-2.9+o0.1i/0.1i+w2c -O -K >> $ps
# Add subtitle
gmt pstext -R0/10/0/15 -JX10/10 -X0.0c -Y4.8c -N -O \
    -F+f10p,25,black+jLB >> $ps << EOF
1.0 13.6 Global gravity grid from CryoSat-2 and Jason-1, 1 min resolution, SIO, NOAA, NGA.
EOF
# Convert to image file using GhostScript
gmt psconvert Grav_VE.ps -A0.5c -E720 -Tj -Z
\end{lstlisting}

\subsection[\appendixname~\thesubsection]{GMT script for vertical gradient of gravity}
\begin{lstlisting}[language=bash, caption=\textbf{GMT code used for plotting Figure \ref{fig6} (vertical gradient of gravity modelling)}]
#!/bin/sh
# Purpose: geophysical mapping of Venezuela (vertical free-air gravity gradient)
# GMT modules: gmtset, gmtdefaults, img2grd, makecpt, grdimage, psscale, grdcontour, psbasemap, gmtlogo, psconvert, pscoast
# http://soliton.vm.bytemark.co.uk/pub/cpt-city/pj/4/index.html
# GMT set up
gmt set FORMAT_GEO_MAP=dddF \
# Overwrite defaults of GMT
gmtdefaults -D > .gmtdefaults
gmt img2grd curv_27.1.img -R286/300.5/0/12.5 -Ggrav_v_VE.grd -T1 -I1 -E -S0.1 -V
gdalinfo grav_v_VE.grd -stats
# Minimum=-488.036, Maximum=870.350, Mean=0.358, StdDev=43.643
# Generate a color palette table from grid
gmt makecpt -Chaxby -T-100/100 > colors.cpt
# Generate a file
ps=Grav_VE_v.ps
gmt grdimage grav_v_VE.grd -Ccolors.cpt -R286/300.5/0/12.5 -JM6.0i -P -I+a15+ne0.75 -Xc -K > $ps
# Add isolines
gmt grdcontour grav_v_VE.grd -R -J -C100 -Wthinnest -O -K >> $ps
# Add grid
gmt psbasemap -R -J \
    --MAP_FRAME_AXES=WEsN \
    --FORMAT_GEO_MAP=ddd:mm:ssF \
    -Bpx4f2a2 -Bpyg4f2a2 -Bsxg4 -Bsyg2 \
    -B+t"Vertical gravity gradient for Venezuela, Trinidad and Tobago" -O -K >> $ps  
# Add legend
gmt psscale -Dg286/-1.0+w15.0c/0.4c+h+o0.0/0i+ml+e -R -J -Ccolors.cpt \
    -Bg25f2.5a25+l"Color scale 'haxby' B. Haxby's color scheme for geoid & gravity [C=RGB -T-200/200]" \
    -I0.2 -By+l"mGal" -O -K >> $ps
# Add scale, directional rose
gmt psbasemap -R -J \
    -Lx14.0c/-2.3c+c50+w200k+l"Mercator projection. Scale (km)"+f \
    -UBL/-5p/-70p -O -K >> $ps
# Add coastlines, borders, rivers
gmt pscoast -R -J -P -Ia/thinnest,blue -Na -N1/thick,white -Wthin,darkslategray -Df -O -K >> $ps
# Texts
gmt pstext -R -J -N -O -K \
-F+jTL+f12p,21,white+jLB+a-315 >> $ps << EOF
286.7 5.8 Los Andes
EOF
# Add GMT logo
gmt logo -Dx7.0/-2.9+o0.1i/0.1i+w2c -O -K >> $ps
# Add subtitle
gmt pstext -R0/10/0/15 -JX10/10 -X0.0c -Y4.8c -N -O \
    -F+f10p,25,black+jLB >> $ps << EOF
1.0 13.6 Global gravity grid from CryoSat-2 and Jason-1, 1 min resolution, SIO, NOAA, NGA.
EOF
# Convert to image file using GhostScript
gmt psconvert Grav_VE_v.ps -A0.5c -E720 -Tj -Z
\end{lstlisting}

\subsection[\appendixname~\thesubsection]{GMT script for seismic mapping}
\begin{lstlisting}[language=bash, caption=\textbf{GMT codes used for plotting Figure \ref{fig7} (seismic mapping)}]
#!/bin/sh
# Purpose: seismicity map of Venezuela (earthquake distribution according to the IRIS database (1997-2021) of seismic events)
# GMT modules: gmtset, gmtdefaults, grdcut, makecpt, grdimage, psscale, grdcontour, psbasemap, gmtlogo, psconvert
# http://soliton.vm.bytemark.co.uk/pub/cpt-city/wkp/shadowxfox/tn/colombia.png.index.html
# GMT set up
gmt set FORMAT_GEO_MAP=dddF 
gmtdefaults -D > .gmtdefaults
# Extract a topographic subset of Venezuela and surroundings:
gmt grdcut ETOPO1_Ice_g_gmt4.grd -R286/300.5/0/12.5 -Gve_relief1.nc
gmt grdcut GEBCO_2019.nc -R286/300.5/0/12.5 -Gve_relief.nc
gdalinfo -stats ve_relief.nc
# Min=-4947.963 Max=5535.625
# Make color palette
gmt makecpt -Ccolombia.cpt > pauline.cpt
gmt makecpt -Cseis -T1.9/7.3/0.5 -Z > steps.cpt
ps=Seis_VE.ps
# Make raster image
gmt grdimage ve_relief.nc -Cpauline.cpt -R286/300.5/0/12.5 -JM6.5i -I+a15+ne0.75 -Xc -P -K > $ps
# Add legend
gmt psscale -Dg286/-1.0+w16.5c/0.4c+h+o0.0/0i+ml+e -R -J -Cpauline.cpt \
    -Bg500f50a500+l"Colormap..." \
    -I0.2 -By+lm -O -K >> $ps
# Add isolines
gmt grdcontour ve_relief1.nc -R -J -C500 -Wthinnest,brown -O -K >> $ps
# Add coastlines, borders, rivers
gmt pscoast -R -J -P \
    -Ia/thinner,blue -Na -N1/thickest,olivedrab1 -Wthin,lightcyan2 -Df -O -K >> $ps 
# Add grid
gmt psbasemap -R -J \
    --MAP_FRAME_AXES=WEsN \
    --FORMAT_GEO_MAP=ddd:mm:ssF \
    -Bpx4f2a2 -Bpyg4f2a2 -Bsxg4 -Bsyg2 \
    -B+t"Seismicity in Venezuela and Trinidad and Tobago: IRIS database (1997-2021)" -O -K >> $ps
# Add scale, directional rose
gmt psbasemap -R -J \
    -Lx14.0c/-3.6c+c50+w200k+l"Mercator projection. Scale: km"+f \
    -UBL/-5p/-110p -O -K >> $ps
# Add earthquake points
# separator in numbers of table: dot (.), not comma ! (British style)
gmt psxy -R -J quakes_VE.ngdc -Wfaint -i4,3,6,6s0.05 -h3 -Scc -Csteps.cpt -O -K >> $ps
# Add geological lines and points
gmt psxy -R -J volcanoes.gmt -St0.4c -Gred -Wthinnest -O -K >> $ps
# fabric and magnetic lineation picks fracture zones
gmt psxy -R -J GSFML_SF_FZ_KM.gmt -Wthicker,goldenrod1 -O -K >> $ps
gmt psxy -R -J GSFML_SF_FZ_RM.gmt -Wthicker,pink -O -K >> $ps
gmt psxy -R -J ridge.gmt -Sf0.5c/0.15c+l+t -Wthick,red -Gyellow -O -K >> $ps
gmt psxy -R -J ridge.gmt -Sc0.05c -Gred -Wthickest,red -O -K >> $ps
# tectonic plates
gmt psxy -R -J TP_Caribbean.txt -L -Wthickest,purple -O -K >> $ps
gmt psxy -R -J TP_Cocos.txt -L -Wthickest,purple -O -K >> $ps
gmt psxy -R -J TP_Nazca.txt -L -Wthickest,purple -O -K >> $ps
gmt psxy -R -J TP_South_Am.txt -L -Wthickest,purple -O -K >> $ps
# Texts
gmt pstext -R -J -N -O -K \
-F+f10p,17,black+jLB+a-0 >> $ps << EOF
293.20 10.58 Caracas
EOF
# Processed likewise for similar texts
gmt pslegend -R -J -Dx1.5/-3.0+w17.8c+o-2.0/0.1c \
    -F+pthin+ithinner+gwhite \
    --FONT=8p,black -O -K << FIN >> $ps
H 10 Helvetica Seismicity: earthquakes magnitude (M) from 1.9 to 7.3
N 9
S 0.3c c 0.3c red 0.01c 0.5c M (7.1-7.3)
S 0.3c c 0.3c tomato 0.01c 0.5c M (6.4-7.0)
S 0.3c c 0.3c orange 0.01c 0.5c M (5.6-6.3)
S 0.3c c 0.3c yellow 0.01c 0.5c M (4.8-5.5)
S 0.3c c 0.3c chartreuse1 0.01c 0.5c M (4.1-4.7)
S 0.3c c 0.3c chartreuse1 0.01c 0.5c M (3.4-4.0)
S 0.3c c 0.3c cyan3 0.01c 0.5c M (2.7-3.3)
S 0.3c c 0.3c blue 0.01c 0.5c M (1.9-2.6)
S 0.3c t 0.3c red 0.03c 0.5c Volcanoes
FIN
# Add GMT logo
gmt logo -Dx7.0/-4.5+o0.1i/0.1i+w2c -O -K >> $ps
# Add subtitle
gmt pstext -R0/10/0/15 -JX10/10 -X0.5c -Y5.9c -N -O \
    -F+f11p,25,black+jLB >> $ps << EOF
1.0 13.6 DEM: SRTM/GEBCO, 15 arc sec grid. Earthquakes: IRIS Seismic Event Database
EOF
# Convert to image file using GhostScript
gmt psconvert Seis_VE.ps -A1.7c -E720 -Tj -Z
\end{lstlisting}

%%%%%%%%%%%%%%%%%%%%%%%%%%%%%%%%%%%s%%%%%%%
\begin{adjustwidth}{-\extralength}{0cm}
%\printendnotes[custom] % Un-comment to print a list of endnotes

\reftitle{References}

% Please provide either the correct journal abbreviation (e.g. according to the “List of Title Word Abbreviations” http://www.issn.org/services/online-services/access-to-the-ltwa/) or the full name of the journal.
% Citations and References in Supplementary files are permitted provided that they also appear in the reference list here. 

%=====================================
% References, variant A: external bibliography
%=====================================
\bibliography{Venezuela.bib}
%\nocite{*}
 
 \section*{Short Biography of Authors}
\bio
{\raisebox{-0.35cm}{\includegraphics[width=3.5cm,height=5.3cm,clip,keepaspectratio]{Photo_author_1.jpg}}}
{\textbf{Polina Lemenkova} is an Engineer and a PhD student of Prof. Dr. O. Debeir in the Image Analysis Research Unit (LISA-IA) in École polytechnique de Bruxelles, Université Libre de Bruxelles. Her research interests are focused on image processing, cartography, spatial data analysis, remote sensing and applied programming. She has a strong research experience in Earth and environmental sciences, mapping and data visualization, specifically in geological and geophysical domains. Her current topic is on seismic data processing using machine learning methods and scripting languages.}
%
\bio
{\raisebox{-0.35cm}{\includegraphics[width=3.5cm,height=5.3cm,clip,keepaspectratio]{Photo_author_2.jpg}}}
{\textbf{Olivier Debeir} is a Head of the Image Analysis Research Unit (LISA-IA) in the École polytechnique de Bruxelles (Brussels Faculty of Engineering), Université Libre de Bruxelles. His research interest is in addressing questions related to artificial intelligence and machine learning in general, and specifically on computer vision, graphics, image analysis, segmentation and classification, pattern recognition and deep learning. His research activity also concerns applied statistics and diverse algorithms of data analysis and modelling in natural sciences. He is an expert in programming languages.}

%%%%%%%%%%%%%%%%%%%%%%%%%%%%%%%%%%%%%%%%%%
\end{adjustwidth}
\end{document}